% problem_id: S001-chow-lemma-tricky
% source_id: S001 (Stacks Project)
% source_file: chow.tex
% statement_locator: chow.tex lines 17107--17139
% proof_locator: chow.tex lines 17141--17354
% env: lemma
% id_note: label 跨文件重名 ⇒ id 用文件限定形式
% pairing: tight (verified)
% status: complete (statement + author proof verbatim + reason + locators)
%
\begin{lemma}
\label{lemma-tricky}
Let $R$ be a local ring with residue field $\kappa$.
Let $\alpha : (M, \varphi, \psi) \to (M', \varphi', \psi')$
be a morphism of $(2, 1)$-periodic complexes over $R$.
Assume
\begin{enumerate}
\item $M$, $M'$ have finite length,
\item $(M, \varphi, \psi)$, $(M', \varphi', \psi')$ are exact,
\item the maps $\varphi$, $\psi$ induce the zero map on
$K = \Ker(\alpha)$, and
\item the maps $\varphi$, $\psi$ induce the zero map on
$Q = \Coker(\alpha)$.
\end{enumerate}
Denote $N = \alpha(M) \subset M'$. We obtain two short exact sequences
of $(2, 1)$-periodic complexes
$$
\begin{matrix}
0 \to (N, \varphi', \psi') \to (M', \varphi', \psi') \to (Q, 0, 0) \to 0 \\
0 \to (K, 0, 0) \to (M, \varphi, \psi) \to (N, \varphi', \psi') \to 0
\end{matrix}
$$
which induce two isomorphisms $\alpha_i : Q \to K$, $i = 0, 1$. Then
$$
\det\nolimits_\kappa(M, \varphi, \psi)
=
\det\nolimits_\kappa(\alpha_0^{-1} \circ \alpha_1)
\det\nolimits_\kappa(M', \varphi', \psi')
$$
In particular, if $\alpha_0 = \alpha_1$, then
$\det\nolimits_\kappa(M, \varphi, \psi) =
\det\nolimits_\kappa(M', \varphi', \psi')$.
\end{lemma}

\begin{proof}
There are (at least) two ways to prove this lemma. One is to produce an
enormous commutative diagram using the properties of the determinants.
The other is to use the characterization of the determinants in terms
of admissible sequences of elements. It is the second approach that we
will use.

\medskip\noindent
First let us explain precisely what the maps $\alpha_i$ are.
Namely, $\alpha_0$ is the composition
$$
\alpha_0 : Q = H^0(Q, 0, 0) \to H^1(N, \varphi', \psi') \to H^2(K, 0, 0) = K
$$
and $\alpha_1$ is the composition
$$
\alpha_1 : Q = H^1(Q, 0, 0) \to H^2(N, \varphi', \psi') \to H^3(K, 0, 0) = K
$$
coming from the boundary maps of the short exact sequences of complexes
displayed in the lemma. The fact that the
complexes $(M, \varphi, \psi)$, $(M', \varphi', \psi')$ are exact
implies these maps are isomorphisms.

\medskip\noindent
We will use the notation $I_\varphi = \Im(\varphi)$,
$K_\varphi = \Ker(\varphi)$ and similarly for the other maps.
Exactness for $M$ and $M'$
means that $K_\varphi = I_\psi$ and three similar equalities.
We introduce $k = \text{length}_R(K)$, $a = \text{length}_R(I_\varphi)$,
$b = \text{length}_R(I_\psi)$. Then we see that $\text{length}_R(M) = a + b$,
and $\text{length}_R(N) = a + b - k$, $\text{length}_R(Q) = k$
and $\text{length}_R(M') = a + b$. The exact sequences below will show
that also $\text{length}_R(I_{\varphi'}) = a$ and
$\text{length}_R(I_{\psi'}) = b$.

\medskip\noindent
The assumption that $K \subset K_\varphi = I_\psi$ means that
$\varphi$ factors through $N$ to give an exact sequence
$$
0 \to \alpha(I_\psi) \to N \xrightarrow{\varphi\alpha^{-1}} I_\psi \to 0.
$$
Here $\varphi\alpha^{-1}(x') = y$ means $x' = \alpha(x)$ and $y = \varphi(x)$.
Similarly, we have
$$
0 \to \alpha(I_\varphi) \to N \xrightarrow{\psi\alpha^{-1}} I_\varphi \to 0.
$$
The assumption that $\psi'$ induces the zero map on
$Q$ means that $I_{\psi'} = K_{\varphi'} \subset N$.
This means the quotient $\varphi'(N) \subset I_{\varphi'}$
is identified with $Q$. Note that $\varphi'(N) = \alpha(I_\varphi)$.
Hence we conclude there is an isomorphism
$$
\varphi' : Q \to I_{\varphi'}/\alpha(I_\varphi)
$$
simply described by
$\varphi'(x' \bmod N) = \varphi'(x') \bmod \alpha(I_\varphi)$.
In exactly the same way we get
$$
\psi' : Q \to I_{\psi'}/\alpha(I_\psi)
$$
Finally, note that $\alpha_0$ is the composition
$$
\xymatrix{
Q \ar[r]^-{\varphi'} &
I_{\varphi'}/\alpha(I_\varphi)
\ar[rrr]^-{\psi\alpha^{-1}|_{I_{\varphi'}/\alpha(I_\varphi)}} & & &
K
}
$$
and similarly
$\alpha_1 = \varphi\alpha^{-1}|_{I_{\psi'}/\alpha(I_\psi)} \circ \psi'$.

\medskip\noindent
To shorten the formulas below we are going to write $\alpha x$ instead
of $\alpha(x)$ in the following. No confusion should result since
all maps are indicated by Greek letters and elements by Roman letters.
We are going to choose
\begin{enumerate}
\item an admissible sequence $z_1, \ldots, z_k \in K$
generating $K$,
\item elements $z'_i \in M$ such that $\varphi z'_i = z_i$,
\item elements $z''_i \in M$ such that $\psi z''_i = z_i$,
\item elements $x_{k + 1}, \ldots, x_a \in I_\varphi$ such
that $z_1, \ldots, z_k, x_{k + 1}, \ldots, x_a$ is an admissible
sequence generating $I_\varphi$,
\item elements $\tilde x_i \in M$ such that $\varphi \tilde x_i = x_i$,
\item elements $y_{k + 1}, \ldots, y_b \in I_\psi$ such that
$z_1, \ldots, z_k, y_{k + 1}, \ldots, y_b$ is an admissible
sequence generating $I_\psi$,
\item elements $\tilde y_i \in M$ such that $\psi \tilde y_i = y_i$, and
\item elements $w_1, \ldots, w_k \in M'$ such that
$w_1 \bmod N, \ldots, w_k \bmod N$ are an admissible sequence
in $Q$ generating $Q$.
\end{enumerate}
By Remark \ref{remark-more-elementary} the element
$D = \det_\kappa(M, \varphi, \psi) \in \kappa^*$ is
characterized by
\begin{eqnarray*}
& &
[z_1, \ldots, z_k,
x_{k + 1}, \ldots, x_a,
z''_1, \ldots, z''_k,
\tilde y_{k + 1}, \ldots, \tilde y_b] \\
& = &
(-1)^{ab} D
[z_1, \ldots, z_k,
y_{k + 1}, \ldots, y_b,
z'_1, \ldots, z'_k,
\tilde x_{k + 1}, \ldots, \tilde x_a]
\end{eqnarray*}
Note that by the discussion above
$\alpha x_{k + 1}, \ldots, \alpha x_a, \varphi w_1, \ldots, \varphi w_k$
is an admissible sequence generating $I_{\varphi'}$ and
$\alpha y_{k + 1}, \ldots, \alpha y_b, \psi w_1, \ldots, \psi w_k$
is an admissible sequence generating $I_{\psi'}$.
Hence by Remark \ref{remark-more-elementary} the element
$D' = \det_\kappa(M', \varphi', \psi') \in \kappa^*$ is
characterized by
\begin{eqnarray*}
& &
[\alpha x_{k + 1}, \ldots, \alpha x_a,
\varphi' w_1, \ldots, \varphi' w_k,
\alpha \tilde y_{k + 1}, \ldots, \alpha \tilde y_b,
w_1, \ldots, w_k]
\\
& = &
(-1)^{ab} D'
[\alpha y_{k + 1}, \ldots, \alpha y_b,
\psi' w_1, \ldots, \psi' w_k,
\alpha \tilde x_{k + 1}, \ldots, \alpha \tilde x_a,
w_1, \ldots, w_k]
\end{eqnarray*}
Note how in the first, resp.\ second displayed formula
the first, resp.\ last $k$ entries of the symbols on both sides
are the same. Hence these formulas are really equivalent to the
equalities
\begin{eqnarray*}
& &
[\alpha x_{k + 1}, \ldots, \alpha x_a,
\alpha z''_1, \ldots, \alpha z''_k,
\alpha \tilde y_{k + 1}, \ldots, \alpha \tilde y_b] \\
& = &
(-1)^{ab} D
[\alpha y_{k + 1}, \ldots, \alpha y_b,
\alpha z'_1, \ldots, \alpha z'_k,
\alpha \tilde x_{k + 1}, \ldots, \alpha \tilde x_a]
\end{eqnarray*}
and
\begin{eqnarray*}
& &
[\alpha x_{k + 1}, \ldots, \alpha x_a,
\varphi' w_1, \ldots, \varphi' w_k,
\alpha \tilde y_{k + 1}, \ldots, \alpha \tilde y_b]
\\
& = &
(-1)^{ab} D'
[\alpha y_{k + 1}, \ldots, \alpha y_b,
\psi' w_1, \ldots, \psi' w_k,
\alpha \tilde x_{k + 1}, \ldots, \alpha \tilde x_a]
\end{eqnarray*}
in $\det_\kappa(N)$. Note that
$\varphi' w_1, \ldots, \varphi' w_k$
and
$\alpha z''_1, \ldots, z''_k$
are admissible sequences generating the module
$I_{\varphi'}/\alpha(I_\varphi)$. Write
$$
[\varphi' w_1, \ldots, \varphi' w_k]
= \lambda_0 [\alpha z''_1, \ldots, \alpha z''_k]
$$
in $\det_\kappa(I_{\varphi'}/\alpha(I_\varphi))$
for some $\lambda_0 \in \kappa^*$. Similarly,
write
$$
[\psi' w_1, \ldots, \psi' w_k]
= \lambda_1 [\alpha z'_1, \ldots, \alpha z'_k]
$$
in $\det_\kappa(I_{\psi'}/\alpha(I_\psi))$
for some $\lambda_1 \in \kappa^*$. On the one hand
it is clear that
$$
\alpha_i([w_1, \ldots, w_k]) = \lambda_i[z_1, \ldots, z_k]
$$
for $i = 0, 1$ by our description of $\alpha_i$ above,
which means that
$$
\det\nolimits_\kappa(\alpha_0^{-1} \circ \alpha_1)
=
\lambda_1/\lambda_0
$$
and
on the other hand it is clear that
\begin{eqnarray*}
& &
\lambda_0 [\alpha x_{k + 1}, \ldots, \alpha x_a,
\alpha z''_1, \ldots, \alpha z''_k,
\alpha \tilde y_{k + 1}, \ldots, \alpha \tilde y_b] \\
& = &
[\alpha x_{k + 1}, \ldots, \alpha x_a,
\varphi' w_1, \ldots, \varphi' w_k,
\alpha \tilde y_{k + 1}, \ldots, \alpha \tilde y_b]
\end{eqnarray*}
and
\begin{eqnarray*}
& &
\lambda_1[\alpha y_{k + 1}, \ldots, \alpha y_b,
\alpha z'_1, \ldots, \alpha z'_k,
\alpha \tilde x_{k + 1}, \ldots, \alpha \tilde x_a] \\
& = &
[\alpha y_{k + 1}, \ldots, \alpha y_b,
\psi' w_1, \ldots, \psi' w_k,
\alpha \tilde x_{k + 1}, \ldots, \alpha \tilde x_a]
\end{eqnarray*}
which imply $\lambda_0 D = \lambda_1 D'$. The lemma follows.
\end{proof}
